% !TeX root = ../教材题库_全部条目_转写工作稿.tex
% Source: XJ-B1-C03-EX-P075-Y0639
% 第3章 / 3.1 函数 / 教材 p.75 / PDF 第82页
% OCR未校对：公式与文字均可能有误；未完成数学TeX转写；禁止直接发给学生
% 原题及相邻原解截图（核对用）：../题目原图/XJ-B1-C03-EX-P075-Y0639-01.webp
% 本文件只在转写完成后作为题干使用；OCR 条目尚不是可靠题干。
\begin{bankitem}{3.1 例题 例题8}{第3章 / 3.1 函数 / 教材 p.75 / PDF 第82页}{XJ-B1-C03-EX-P075-Y0639}{OCR未校对：公式与文字均可能有误；未完成数学TeX转写；禁止直接发给学生}
\begin{ocrdraft}
[教材p.75 / PDF第82页]
\par
0二0]画出函数一一2十十的图象．
\par
解为了去掉绝对值符号，需分段讨论：
\par
C（5，9）
\par
当<一1时，丆（工）一（2一工）十（一工一1)一1一2；
\par
韭00到0的
\par
当一1主2时，不工）一2一工十十1一3；
\par
召（2，3）
\par
当x>2时，丿飞r）一x一2十x十1一2一1.
\par
0
\par
分段画出不x）的图象，如图3．1一4所示．
\par
图3．1一4
\par
作分段函数的图象，一段一段地按表达式画是基本方
\par
法．本例可更简单：因为@一2]和@十0不论在哪一段中，去掉绝对值符号后所得
\par
的都是一次函数，其和（或差）是一次函数或常数函数，图象都是直线，每段都是
\par
线段或射线，要作出其图象，只要在每段上取两个点就够了．
\par
为了简单，尽可能取分段点．
\par
绝对值函数的分段点，就是绝对值符号中的量为
\par
0的点，所以不到的分段点在一一1和工一2之处．为了画出第一段和第三段的射
\par
线，再取工一一4和一5．这样，即使不去掉绝对值符号，也能画图．
\par
知直一
\par
常用的分段函数
\par
在前面，我们介绍了不一0《是一个分段函数0到也写作abs(x)）．在数学里，还
\par
有几个很常用的分段函数．例如，常用[到表示不大于的最大整数，如匚4]：4，[4．1」一
\par
4，[一3．14」一一4等，匚到叫作整数部分函数．对应地，\{到：一匚到叫作小数部分函
\par
数，例如\{的一0，\{4．I\}=0．1，\{一3．14\}=0．86等．还有一个符号函数sgn(x)，当>0
\par
时，sgn(x)=I；当=0时，sgn(x)=0；当<0时，sgn(x)=
\end{ocrdraft}
\end{bankitem}
